% ===================================================================
%  تابلوی شمارهٔ ۱۱ — شهر و بناهای آن. --- آتش‌سوزی.
%
%  Traduction, faite en 2026, en persan d'aujourd'hui, sur l'ido de
%  texto/io. Regles de la colonne : voir 10-tablo-01.tex. Deux
%  scenes, deux numerotations. Ouverture de la TROISIEME SERIE :
%  « سری سوم ». Deux blocs marques « ¶ » par ordo.py.
%
%  C'EST LE TABLEAU QUI N'AVAIT AUCUN GROS PLAN DANS LES 42 COLONNES,
%  ET IL EN A MAINTENANT 89. La planche 11 numerote c1/c2 pour ses
%  deux scenes mais grave ses numeros NUS ; html.py cherchait « c1:7 »
%  et ne trouvait rien, dans toutes les langues a la fois. On ne s'en
%  est apercu qu'en comptant les boutons tableau par tableau, pendant
%  l'ecriture de la colonne ourdoue.
%
%  LA VILLE PERSANE A DEUX AGES, ET LE TABLEAU LES SEPARE EXACTEMENT
%  COMME LE TABLEAU 11 JAVANAIS SEPARAIT LES SIENS. Ce qui existait
%  avant le XIXe siecle porte un nom d'ici : بازار, میدان la place
%  (96), کوچه la ruelle (11), محله le quartier, دروازه la porte de
%  ville (8), برج la tour (9), حصار l'enceinte, کاروانسرا, حمام,
%  گورستان. Ce que l'Etat moderne a bati porte un nom francais ou
%  arabe : بورس la Bourse (14), موزه le musee, پست la poste (43), بانک
%  la banque (32), بلوار le boulevard (19), تراموا le tram (31),
%  کازرم la caserne (16), گمرک la douane (10).
%
%  ET گمرک EST L'EXCEPTION QUI SE REPETE. Au tableau 11 javanais,
%  « pabean » etait le seul batiment moderne a garder un nom d'ici,
%  parce que la douane existait aux ports de Java avant les
%  Hollandais. Le persan fait exactement la meme exception, avec un
%  mot different : گمرک vient du grec kommerkion par l'arabe, et il
%  est en persan depuis mille ans. Deux colonnes, deux mots, une seule
%  raison — la douane est le plus vieux des batiments modernes.
%
%  L'INCENDIE RETOURNE AU PERSAN D'UN SEUL COUP, COMME EN JAVANAIS :
%  آتش le feu, دود la fumee (75), شعله la flamme, نردبان l'echelle
%  (72), آب‌پاش la lance, سوختن bruler. Une ville brule dans la langue
%  de ceux qui l'eteignent, meme quand ses batiments portent des noms
%  etrangers. Troisieme fois que ce partage se verifie — le pawon et
%  le kompor au tableau 6 javanais, la voile et la vapeur au 10, le
%  feu et la Bourse ici.
%
%  ET LE THEATRE (73) EST LE SECOND REFUS DE LA COLONNE, APRES LE
%  SAMOVAR DU TABLEAU 6. Le persan a تعزیه, le drame religieux joue
%  dans les takyeh depuis quatre siecles, et il a سیاه‌بازی, la comedie
%  populaire ; la planche grave un opera a l'italienne avec fosse
%  d'orchestre, loges et peristyle. La colonne ecrit تئاتر, mot
%  francais, et اپرا. Le mot d'ici ne se pose pas sur une chose de
%  la-bas — meme regle qu'au wayang du tableau 16 javanais, et le
%  releve la voit maintenant pour la septieme fois.
% ===================================================================

%%K t11-apar-1 apar
\VUsaut{2.0mm}
\VUcentre{13.2pt}{90}{سری سوم}
\VUsaut{3.0mm}
\VUfilet{16mm}
\VUsaut{5.6mm}
\VUtitre{15.0pt}{60}{تابلوی شمارهٔ ۱۱}
\VUsaut{1.4mm}
\VUpk{11.4pt}{40}{شهر و بناهای آن. --- آتش‌سوزی.}
\VUsaut{2.2mm}

%%K t11-tit-1 sub
\VUsaut{3.0mm}
\VUpk{10.2pt}{40}{حومهٔ شهر.}
\VUsaut{3.0mm}

%%K t11-c1-01-1 p
1. --- من در شهر بزرگی می‌نشینم که صد کیلومتر از اقیانوس دور است و با
این همه \VUgras{بندر} \textsuperscript{(1)} درجه یکی است. رودی که آن را
حاشیه می‌زند چهارصد متر پهنا دارد و خم بسیار زیبایی می‌سازد.

%%K t11-c1-02-1 p
2. --- همهٔ آن بخش شهر که بر \VUgras{اسکله‌ها} \textsuperscript{(2)}
است، سراسر دریایی و صنعتی می‌نماید و \VUgras{حومه‌ای}
\textsuperscript{(3)} می‌سازد که تنها دریانوردان و کارگران در آن
می‌نشینند. اینان در \VUgras{کارخانه‌ها} \textsuperscript{(4)} و در
\VUgras{کارخانهٔ گاز} \textsuperscript{(5)} کار می‌کنند که
\VUgras{دودکش} \textsuperscript{(6)} بلند و \VUgras{گازسنج} آن از بسیار
دور دیده می‌شود.

%%K t11-c1-03-1 p
3. --- در سده‌های میانه، شهر دژبندی شده بود، و هنوز چند بازمانده از
باروهای آن می‌یابیم، به ویژه \VUgras{دروازهٔ} \textsuperscript{(8)}
\VUgras{کاخ دژبندی‌شدهٔ} \textsuperscript{(7)} کهن که دو \VUgras{برج}
\textsuperscript{(9)} در دو سویش هستند و اکنون \VUgras{زندان}
شده‌اند.

%%K t11-c1-04-1 p
4. --- در همهٔ این \VUgras{محله} که بسیار کهن می‌نماید، تنها یک ساختمان
نسبتاً نو است : آن \VUgras{گمرک} \textsuperscript{(10)} است. میان اسکله
و \VUgras{کوچه‌ای} \textsuperscript{(11)} جای دارد که به \VUgras{تالار
شهر} \textsuperscript{(12)} کهن می‌رسد. آن بنای واپسین را به آسانی از
\VUgras{ناقوس‌خانهٔ} \textsuperscript{(13)} غول‌آسایش بازمی‌شناسیم که
بر شهر کهن چیره است.

%%K t11-c1-05-1 p
5. --- در سوی دیگر خیابان، ساختمانی \VUgras{ایستاده} است که به همان
سبک گمرک ساخته شده : آن \VUgras{بورس} \textsuperscript{(14)} است. یکی
از نماهای آن به میدان کوچکی می‌نگرد که \VUgras{فرمانداری}
\textsuperscript{(15)} در آن است. در واقع شهر ما بی‌آنکه پایتخت بزرگی
باشد، با این همه مرکز مهم استان است.

%%K t11-tit-2 sub
\VUsaut{5.0mm}
\VUpk{10.2pt}{40}{بلندترین بخش شهر.}
\VUsaut{3.4mm}

%%K t11-c1-06-1 p
6. --- پادگان که شمارش بس زیاد است، در \VUgras{سربازخانه‌های}
\textsuperscript{(16)} پهناور و بسیار سالم جای دارد ؛ آنها تا نردهٔ
\VUgras{پارک شهر} \textsuperscript{(17)} گسترده‌اند. \VUgras{بلوار}
\textsuperscript{(19)} که به این پارک می‌رسد، از پشت \VUgras{صندوق
پس‌انداز} \textsuperscript{(18)} می‌گذرد و به میدان \VUgras{کلیسای
جامع} \textsuperscript{(20)} می‌رسد.

%%K t11-c1-07-1 p
7. --- همهٔ این بناها مانند آمفی‌تئاتر بر تپهٔ کوچکی گسترده‌اند که
\VUgras{دبیرستان} \textsuperscript{(21)} پهناور ما نیز در آن است.

%%K t11-c1-07-2 p
اگر از راهی اندکی سرازیر پایین بروند که به خیابان بزرگ می‌پیوندد، به
\VUgras{دادگستری} \textsuperscript{(22)} می‌رسند که در برابر آن
\VUgras{باغ همگانی} \textsuperscript{(26)} دلپذیری گسترده است. آن
\VUgras{میدان‌باغ} چندان بزرگ نیست، اما در میان شهر واحه‌ای از سبزی و
خنکی می‌سازد. از این رو شهرنشینان بسیار به آن می‌روند که دوست دارند
زیر سایبان \VUgras{قهوه‌خانه} \textsuperscript{(25)} بنشینند تا به
نوازندگان سرباز گوش دهند که پنجشنبه و یکشنبه در \VUgras{آلاچیق}
\textsuperscript{(27)} می‌نوازند که میان چمن‌ها و بوته‌ها گذاشته شده
است.

%%K t11-c1-08-1 p
8. --- بر یکی از آن چمن‌ها \VUgras{مجسمهٔ} \textsuperscript{(28)}
برنزی ایستاده است، بنایی که به یاد یکی از هم‌شهریان ما برپا شده است که
به زادگاه خود خدمت‌های بزرگ کرد.

%%K t11-tit-3 sub
\VUsaut{4.0mm}
\VUpk{10.2pt}{40}{وسایل نقلیه.}
\VUsaut{3.4mm}

%%K t11-c1-09-1 p
9. --- تا کنون شهر ما هنوز با روشنایی برق روشن نمی‌شود ؛ تنها
\VUgras{چراغ‌های گاز} \textsuperscript{(29)} دارد. اما
\VUgras{روزنامه‌های اینجا} برای بهتر شدن این شیوهٔ روشنایی مبارزه
می‌کنند، \VUgras{چنانکه} از پیش \VUgras{شیوهٔ کشش تراموا‌ها} را
کامل کرده‌اند. \VUgras{گردونه‌های} کهن که \VUgras{اسب‌ها} آنها را
\VUgras{می‌کشیدند}، عملاً با \VUgras{تراموای برقی}
\textsuperscript{(31)} جایگزین شده‌اند که مسافران را با بهایی بسیار
\VUgras{ارزان} تند از یک سر \VUgras{شهر به سر دیگر} می‌برد.

%%K t11-c1-10-1 p
10. --- نزدیک دادگستری \VUgras{بانک} \textsuperscript{(32)} است که در
آن اسناد بازرگانی را تنزیل می‌کنند. \VUgras{مأموران وصول بانک}
\textsuperscript{(33)} می‌روند تا بدهی‌ها را در خانه بگیرند.

%%K t11-tit-4 sub
\VUsaut{4.0mm}
\VUpk{10.2pt}{40}{هنر و آموزش.}
\VUsaut{3.4mm}

%%K t11-c1-11-1 p
11. --- ساختمان زیبایی که در نزدیکی ایستاده است، \VUgras{موزه} است.
آن با هم \VUgras{کتابخانهٔ} \textsuperscript{(34)} شهر،
\VUgras{مجموعه‌های علوم طبیعی} \textsuperscript{(35)} زیبا و
\VUgras{موزهٔ نقاشی} \textsuperscript{(36)} خوش‌دوختی را در بر دارد که
\VUgras{نمایشگاه} دائمی شکوهمندی می‌سازد.

%%K t11-c1-11-2 p
هفته‌ای دو بار برای همگان گشوده می‌شود، اما بیگانگان می‌توانند هر روز
با انعامی به \VUgras{نگهبان} \textsuperscript{(37)} از آن دیدن کنند.
\VUgras{نقاشان} \textsuperscript{(38)} به آنجا می‌آیند تا کار کنند.
آنان را \VUgras{در برابر} سه‌پایهٔ خود می‌بینیم، با \VUgras{تخته‌رنگ}
در دست، که نقاشی‌های نامدار را رونویسی می‌کنند.

%%K t11-c1-12-1 p
12. --- بر بلندی، پشت موزه، \VUgras{گنبد} \textsuperscript{(40)}
\VUgras{بیمارستان} \textsuperscript{(39)} و ساختمان‌های \VUgras{دانشسرا}
\textsuperscript{(41)} را می‌بینیم که آموزگاران مدرسه‌های شهرداری ما در
آن پرورش می‌یافتند.

بر میدان تئاتر \VUgras{ادارهٔ پست} \textsuperscript{(43)} هست که در
\VUgras{خانهٔ خصوصی} \textsuperscript{(42)} جای گرفته است.

%%K t11-tit-5 sub
\VUsaut{5.0mm}
\VUpk{11.4pt}{40}{آتش‌سوزی.}
\VUsaut{3.0mm}

%%K t11-tit-6 sub
\VUsaut{3.4mm}
\VUpk{10.2pt}{40}{فریاد آتش.}
\VUsaut{3.4mm}

%%K t11-c2-01-1 p
1. --- خبر هراس‌انگیزی در شهر پخش می‌شود : همین حالا در تئاتر
آتش‌سوزی روی داده است! درست هنگامی که به \VUgras{میدان}
\textsuperscript{(96)} می‌رسم، جمعیت از پیش بر \VUgras{پیاده‌رو}
\textsuperscript{(46)} و بر \VUgras{سواره‌رو} \textsuperscript{(47)}
گرد آمده است. \VUgras{کارمندان پست} \textsuperscript{(44)} و
\VUgras{نامه‌رسانان} \textsuperscript{(45)} از ادارهٔ پست بیرون
می‌آیند. \VUgras{راهبر تراموا} \textsuperscript{(48)} ناچار شد
گردونه‌اش را نگه دارد تا \VUgras{آمبولانس} \textsuperscript{(50)} شهر
بگذرد که با \VUgras{جراح} \textsuperscript{(52)} و \VUgras{پرستار}
\textsuperscript{(51)} می‌رسد تا \VUgras{زخمی} \textsuperscript{(53)} را
درمان کنند که با \VUgras{برانکارد} \textsuperscript{(54)} نزدیک
\VUgras{کیوسک روزنامه} \textsuperscript{(55)} آورده شده است.

%%K t11-c2-02-1 p
2. --- گروهانی از پیاده‌نظام که از تمرین بازمی‌گشت، ایستاد تا با
\VUgras{ژاندارم‌های سوارهٔ شهری} \textsuperscript{(61)} نظم را
برقرار کند. \VUgras{سواران} که اسبانشان روی دو پا بلند می‌شوند، بهتر
از \VUgras{سربازان} \textsuperscript{(57)} یا \VUgras{ژاندارم‌ها}
\textsuperscript{(63)} \VUgras{کنجکاوان} \textsuperscript{(56)} را عقب
می‌رانند. وانگهی ژاندارم‌ها بسیار سرگرم‌اند. یکی از آنان به
\VUgras{پاسبان‌ها} \textsuperscript{(64)} نشان می‌دهد که
\VUgras{مصدوم} \textsuperscript{(65)} دیگری را کجا باید برد، آتش‌نشان
دلیری که از نردبان افتاده است.

%%K t11-c2-03-1 p
3. --- \VUgras{افسر} \textsuperscript{(66)} جای گذاشتن \VUgras{تلمبهٔ
بخار} \textsuperscript{(67)} را که می‌رسد دقیق می‌کند. در انتظار آن
ماشین نیرومند، بی‌درنگ \VUgras{تلمبهٔ دستی} \textsuperscript{(68)} را
کار گذاشته است که \VUgras{شلنگ} \textsuperscript{(69)} بلندش آب را
سیل‌آسا بر آتش می‌ریزد.

شش آتش‌نشان آن را می‌گردانند، در حالی که رفیقشان که \VUgras{سرلوله}
\textsuperscript{(70)} را نگه داشته است، \VUgras{فوارهٔ آب}
\textsuperscript{(71)} را به مرکز آتش‌سوزی نشانه می‌رود.

%%K t11-c2-04-1 p
4. --- در سوی دیگر تئاتر، \VUgras{نردبان} \textsuperscript{(72)} بزرگ
را تکیه داده‌اند. آن به آتش‌نشانان امکان می‌دهد به شعله‌هایی نزدیک
شوند که میان گرداب‌های \VUgras{دود} \textsuperscript{(75)} سیاه بیرون
می‌زنند.

%%K t11-tit-7 sub
\VUsaut{5.0mm}
\VUpk{10.2pt}{40}{کارهای دلیرانه.}
\VUsaut{3.4mm}

%%K t11-c2-05-1 p
5. --- \VUgras{آتش‌سوزی} \textsuperscript{(74)} در سرسرای
\VUgras{تئاتر} \textsuperscript{(73)} زاده شد، هنگامی که
\VUgras{اپرایی} را تمرین می‌کردند که \VUgras{نمایش} نخستش فردا برگزار
می‌شد. خوشبختانه تنها چند \VUgras{تماشاگر} \textsuperscript{(76)} در
تالار نمایش بودند. آن بینوایان به \VUgras{بالکن} \textsuperscript{(77)}
پناه برده‌اند، جایی که \VUgras{آتش‌نشانان} \textsuperscript{(86)}
دلیر با نردبان و طناب‌های ستبر به یاری آنان می‌آیند.

%%K t11-c2-06-1 p
6. --- زیر \VUgras{ستون‌گاه} \textsuperscript{(78)}، جایی که
\VUgras{اعلان} \textsuperscript{(79)} از نمایش خبر می‌دهد،
\VUgras{نوازندگان} \textsuperscript{(81)} \VUgras{ارکستر}
\textsuperscript{(80)} را می‌بینیم که می‌گریزند و سازهای خود را با خود
می‌برند : \VUgras{ویولن} \textsuperscript{(81)}، \VUgras{فلوت}
\textsuperscript{(82)}، \VUgras{ویولن‌سل} \textsuperscript{(83)}،
\VUgras{ترومبون} \textsuperscript{(84)}، \VUgras{طبل}
\textsuperscript{(85)}، و مانند آن. می‌کوشند بخشی از \VUgras{اثاث}
\textsuperscript{(87)} را نجات دهند.

%%K t11-c2-07-1 p
7. --- \VUgras{هنرپیشگان مرد} \textsuperscript{(89)} و
\VUgras{هنرپیشگان زن} \textsuperscript{(88)}، \VUgras{خوانندگان مرد}
و \VUgras{خوانندگان زن} ناچار شدند با \VUgras{جامهٔ}
\textsuperscript{(90)} نمایش خود بگریزند. تنور به \VUgras{روزنامه‌نگاری}
\textsuperscript{(92)} توضیح می‌دهد که چگونه روی داد. \VUgras{خبرنگار}
بی‌درنگ از نخستین لحظه با مقامات دویده بود : \VUgras{فرماندار}
\textsuperscript{(91)}، \VUgras{شهردار} \textsuperscript{(93)}،
\VUgras{کمیسر پلیس} \textsuperscript{(94)} و \VUgras{مأمور قضایی}
\textsuperscript{(95)} که بازجویی خواهند کرد تا در بارهٔ مسئولیت‌ها
داوری کنند.

\VUsaut{6.0mm}
\VUfilet{22mm}
